\section{Pricing the fiscal gap, 1980--2025}\label{sec:pricing}

\subsection{Parameters}\label{sec:parameters}

The debt sensitivity of rates is $\psi=3$bp per percentage point of debt/GDP, with a range of 2--4.5. The sources are \citet{laubach2009} (3--4), \citet{engen2005} (about 3), \citet{plante2025} (3.0--3.5), \citet{bhatt2026} (3--4) and \citet{bi2026} (2.8 at peak, from identified supply shocks); CBO's 2bp \citep{neveu2024} is the low case. $\psi$ may differ across monetary regimes, for example when the Fed was buying Treasuries, and Table~\ref{tab:robust} lets it do so.

The fiscal response $\hat\varphi$ is a policy choice, and results are reported as a function of it. As a reference path we use \citeauthor{auerbach2024}'s (\citeyear{auerbach2024}) estimates of the legislated response to lagged net interest: 0.39 through 2003 (s.e.\ 0.13) and 0 from 2004. Their post-2004 coefficient, $-0.31$ with s.e.\ 0.41, cannot reject 0.39. We also report a single $\hat\varphi$ common to all years, from 0 to 0.39; \citet{degroot2015} imply about 0.7 over ten years for European countries, which we treat as an upper bound. Our own evidence (Appendix~\ref{app:fiscal}) is that the U.S. response changed, but in the mid-1990s and around 2009--13 rather than at 2004, and that UK fiscal events put it at about 0.15.

The marginal cost $r$ is the steady-state cost of the existing structure: each security's original tenor priced at current yields. Before 2003 we use remaining-maturity pricing, corrected by its measured 0.28pp gap on 2003--2025. Expected nominal growth $g$ is proxied by realized forward ten-year nominal GDP growth through 2015; from 2016 we use the December FOMC Summary of Economic Projections, the longer-run median real GDP growth (1.8--1.9\%) plus the 2\% inflation objective. Trailing ten-year growth and an ex-ante measure (trailing real growth plus the Cleveland Fed's ten-year expected inflation) are sensitivities. In 2013--15, when both measures exist, the SEP-based $g$ (4.0--4.3\%) was below realized forward growth (5.2--5.4\%), which was inflated by 2021--22.

\subsection{The threshold and the price}\label{sec:threshold}

Figure~\ref{fig:limitmap} and Table~\ref{tab:limitperiod} summarize the two layers from 1980 to 2025; Appendix~\ref{app:annual} reports every year. Three readings follow.

\begin{figure}[htbp]
\centering
\includegraphics[width=\textwidth]{fig3_limit_map.pdf}
\caption{Two-layer limit, 1980--2025}\label{fig:limitmap}
\begin{minipage}{0.95\textwidth}\footnotesize Top: fiscal threshold $\varphi^*$ (line; band spans $g$ and $\psi$) and the reference path of the fiscal response $\hat\varphi$ (dashed). Middle: inflation-layer ratio $R=\int P/\int E$, consolidated (erosion base includes currency, and reserves before 2008) and for the Treasury's own marketable debt. Bottom: inflation needed to cover the gap $(\varphi^*-\hat\varphi)^+$ after a $+1$pp permanent rate rise, pp per year for 10 years; bars use the reference path of $\hat\varphi$ (0.39 to 2003, 0 from 2004), the line a common $\hat\varphi=0.25$ in every year.\end{minipage}
\end{figure}

\begin{table}[htbp]
\centering
\caption{Two-layer limit by period (consolidated; g forward; $\psi$ = 3bp)}\label{tab:limitperiod}
\footnotesize
\begin{tabularx}{\textwidth}{l>{\centering\arraybackslash}X>{\centering\arraybackslash}X>{\centering\arraybackslash}X>{\centering\arraybackslash}X>{\centering\arraybackslash}X}
\toprule
Period & $\varphi^*$ & $\hat\varphi$ & Gap & Inflation-layer ratio & Inflation to cover gap (pp/yr) \\
\midrule
1981–85 & 0.41–0.56 & 0.39 & 0.02–0.17 & 1.4–1.6 & 0.04–0.23 \\
1986–99 & 0.12–0.43 & 0.39 & $\approx$ 0 & 1.3–1.6 & $\approx$ 0 \\
2005–07 & 0.28–0.44 & 0 & 0.28–0.44 & 1.4–1.5 & 0.38–0.61 \\
2009–21 & mostly < 0 (2017–19: 0.11–0.21) & 0 & $\approx$ 0 (2017–19: 0.11–0.21) & 2.0–2.6 & $\approx$ 0 (2017–19: 0.23–0.43) \\
2023–25 & 0.46–0.48 & 0 & 0.46–0.48 & 2.2–2.4 & 1.01–1.12 \\
\bottomrule
\end{tabularx}
\end{table}


First, the fiscal threshold is ordinary today. The period that stands out is 2008--21, when $r+\psi b<g$ and debt stabilized with no fiscal response, and that period is over.

Second, the fiscal response. On the reference path the offset roughly met the threshold until the mid-2000s and has fallen short since. The gap opened in 2005--07, when $\varphi^*$ was 0.28--0.44, and the low-rate period of 2008--21 hid it. Because the level and timing of the response are uncertain (Appendix~\ref{app:fiscal}), Figure~\ref{fig:limitmap} also shows a common $\hat\varphi=0.25$ in every year.

Third, QE raised the price of the inflation route. The consolidated ratio rose from 1.2--1.6 in 1980--2007 to 2.0--2.6 from 2008 on. Before QE, consolidation roughly halved the ratio relative to the Treasury's own liabilities (0.49--0.53 times the Treasury-only value in 2003--07), because the Fed held Treasuries against currency, which inflation erodes. Since 2009 the two have been about equal (0.86--1.19), because the Fed's additional Treasuries and MBS are financed by interest-bearing reserves that reprice overnight. The Fed's share of marketable Treasuries was in fact lower at end-2025 (14\%) than at end-2007 (16\%). What QE changed is how the Fed's book is funded: SOMA Treasuries were 0.88 times currency in 2007 and 1.73 times in 2025, and reserves also fund the MBS.

\paragraph{Decomposition, end-2007 to end-2025.} The consolidated ratio rose from 1.45 to 2.18. A chain of counterfactual balance sheets splits the 0.73 rise into three parts. Currency fell from 22\% to 8\% of interest-bearing liabilities, which adds 0.58. The Treasury's maturity and TIPS choices between 2007 and 2025, with a Fed funded only by currency, subtract 0.23: the Treasury's lengthening worked against the rise. QE, in the sense that reserves rather than currency fund the rest of the Fed's book, adds 0.38. The QE share depends on the order: measured first, on the end-2007 balance sheet (the Fed's Treasuries scaled to the end-2025 ratio to currency, and reserves added for the MBS), QE adds 0.57. QE therefore accounts for 52--78\% of the rise; the fall of currency relative to debt, which reflects deficits more than monetary policy, accounts for most of the rest.

The combined requirement in 2023--25 is the largest since 1980, about 1.8 times its 2006--07 level. The equivalent one-time surprise rise in the price level is about 3.3\% (3.27--3.42\% in 2023--25). As Appendix~\ref{app:whichinflation} explains, it is much less sensitive to consolidation: since 2009 the consolidated value has been 0.95--1.05 times the Treasury-only value.

\subsection{When the central bank responds}\label{sec:kappa}

The requirements above assume a central bank that tolerates the inflation. Table~\ref{tab:kappa} prices the same gaps when it responds, using Proposition~\ref{prop:kappa}.

\begin{table}[htbp]
\centering
\caption{Monetary tolerance threshold $\kappa^*$ = 1/R and the requirement under a monetary response}\label{tab:kappa}
{\small consolidated; H = 10; +1pp permanent rate rise; $\hat\varphi$ on the reference path\par}\smallskip
\footnotesize
\begin{tabularx}{\textwidth}{l>{\centering\arraybackslash}X>{\centering\arraybackslash}X>{\centering\arraybackslash}X>{\centering\arraybackslash}X>{\centering\arraybackslash}X>{\centering\arraybackslash}X}
\toprule
Year & Ratio R & $\kappa^*$: H = 10 (H = 5; H = 15) & Inflation to cover gap (pp/yr): $\kappa$ = 0 & $\kappa$ = 0.25 & $\kappa$ = 0.5 (Taylor rule) & $\kappa$ = 0.5, H = 5 \\
\midrule
1981 & 1.39 & 0.72 (1.07; 0.59) & 0.23 & 0.36 & 0.77 & 0.30 \\
1984 & 1.56 & 0.64 (1.03; 0.50) & 0.20 & 0.33 & 0.91 & 0.24 \\
2007 & 1.45 & 0.69 (1.13; 0.53) & 0.61 & 0.96 & 2.23 & 0.67 \\
2019 & 2.01 & 0.50 (0.88; 0.37) & 0.23 & 0.46 & none & 0.30 \\
2023 & 2.38 & 0.42 (0.69; 0.32) & 1.12 & 2.76 & none & 2.53 \\
2024 & 2.25 & 0.44 (0.74; 0.34) & 1.08 & 2.46 & none & 1.97 \\
2025 & 2.18 & 0.46 (0.77; 0.35) & 1.01 & 2.22 & none & 1.70 \\
\bottomrule
\end{tabularx}
\par\smallskip{\footnotesize\raggedright \emph{Notes.} ``None'': $\kappa$ $\ge$ $\kappa^*$, so no finite rate of inflation covers the gap. $\kappa$ is the rise in the real rate per point of sustained surprise inflation; a Taylor (1993) rule has $\kappa$ = 0.5.\par}
\end{table}


From 1980 to 2007 the U.S. monetary tolerance threshold was 0.64--0.82 over a ten-year horizon. In every year since 2009 it has been 0.39--0.50. The threshold rises as the horizon shortens, because less of the debt has repriced: over five years it was 1.01--1.26 before 2008 and 0.66--0.88 since 2009; over fifteen years, 0.49--0.65 and 0.29--0.37. At every horizon QE lowered it by about two fifths, by 37\%, 39\% and 40\% from 2003--07 to 2023--25 at five, ten and fifteen years. Whether it falls below a Taylor rule does depend on the horizon: it does at ten and fifteen years, and not at five, where a Taylor-rule response instead raises the 2023--25 requirement to 1.7--2.5pp a year, against 0.67 in 2007.

QE thus moved the consolidated balance sheet from a position where the inflation route survived a Taylor-rule response over a decade to one where it does not. Before QE the Fed's currency funding kept $\kappa^*$ high: the Treasury-only threshold in 2003--07 was 0.36--0.37, and consolidation doubled it. Since QE consolidation no longer does. Even at $\kappa=0.25$, half a Taylor rule, the 2023--25 requirement more than doubles, to 2.2--2.8pp a year. Beyond a few years, the inflation route to fiscal relief is open only to a central bank that responds to inflation much less than its usual rule.

\subsection{Robustness}\label{sec:robust}

Table~\ref{tab:robust} recomputes the 2023--25 requirement under alternative parameters, holding the measured clock fixed.

\begin{table}[htbp]
\centering
\caption{Robustness of the 2023–25 requirement}\label{tab:robust}
{\small +1pp permanent rate rise, H = 10\par}\smallskip
\small
\begin{tabularx}{\textwidth}{>{\raggedright\arraybackslash}p{0.3\textwidth}>{\centering\arraybackslash}X>{\centering\arraybackslash}X>{\centering\arraybackslash}X>{\centering\arraybackslash}X}
\toprule
Variant & $\varphi^*$, 2023–25 & Inflation to cover gap, 2023–25 (pp/yr) & Highest before 2022 (year) & 2023–25 highest since 1980? \\
\midrule
Baseline (forward/SEP g; $\psi$ = 3bp on consolidated debt) & 0.46–0.48 & 1.01–1.12 & 0.61 (2007) & yes \\
g ex ante: trailing real growth + 10-year expected inflation & 0.32–0.34 & 0.71–0.80 & 0.41 (2018) & yes \\
g: trailing 10-year nominal growth & 0.24–0.28 & 0.52–0.67 & 0.63 (2018) & yes \\
$\psi$ on privately held marketable debt (no reserves) & 0.43–0.45 & 0.95–1.03 & 0.61 (2007) & yes \\
$\psi$ on debt held by the public & 0.47–0.49 & 1.03–1.14 & 0.65 (2007) & yes \\
$\psi$ = 2bp & 0.38–0.40 & 0.83–0.93 & 0.57 (2007) & yes \\
$\psi$ = 4.5bp & 0.55–0.56 & 1.21–1.33 & 0.73 (2018) & yes \\
Monte Carlo: $\psi$ ~ U[2, 4.5]bp, g $\pm$ N(0, 0.5pp), $\hat\varphi$ ~ U[0.25, 0.5] to 2003 and U[0, 0.25] after & — & 0.36–1.27 (90\% band) & — & in 99\% of draws \\
$\psi$ by regime (1980–87, 1988–2007, 2008–21, 2022–25), each 1, 2, 3 or 4.5bp: 256 combinations & 0.28–0.56 & 0.64–1.27 & — & in 94\% (not if $\psi$ is 4.5bp in 2008–21 and 1bp today) \\
Same, one common $\hat\varphi$ = 0 & 0.28–0.56 & 0.64–1.27 & — & in 75\% (not if today's $\psi$ is 1bp) \\
Same, one common $\hat\varphi$ = 0.25 & 0.28–0.56 & 0.07–0.70 & — & in 50\% (not if today's $\psi$ is 2bp or less) \\
\bottomrule
\end{tabularx}
\end{table}


With one $\psi$ for all years, the ranking holds in every variant. With $\psi$ free to differ across four monetary and fiscal regimes, it depends almost only on today's $\psi$: it holds whenever today's $\psi$ is at least 2bp (3bp at a common $\hat\varphi$ of 0.25), whatever the $\psi$ of earlier regimes. The $\psi$ of the 1980s hardly matters, because debt was small. The data cannot narrow $\psi$ by regime: Laubach-style regressions of the five-year-forward five-year Treasury rate on CBO's projected debt, in changes between baselines, give 2.6bp (s.e.\ 2.1) pooled and are uninformative by regime. The threshold in 2023--25 stays within its 1980--2007 range in 88\% of the combinations; it is above that range if today's $\psi$ is 4.5bp.

Expected growth matters most for the level: an ex-ante measure built from trailing real growth and the Cleveland Fed's ten-year expected inflation puts expected nominal growth at 4.8\% in 2025 rather than 3.8\%, and lowers the requirement to 0.71--0.80pp. Where $\psi$ applies barely matters: the elasticities in the literature are estimated on debt held by the public, and reserves are zero-duration safe assets, but applying $\psi$ only to privately held marketable debt moves $\varphi^*$ by about 3 points. In the Monte Carlo the 90\% band for 2023--25 is 0.36--1.27pp a year, and 2023--25 is the highest since 1980 in 99\% of draws. With $H=5$ the 2023--25 requirement is 0.60--0.69pp a year, and with $H=15$ it is 1.33--1.46. If all debt counted as erodible and currency were ignored, the requirement would be 1.22--1.41; ignoring currency alone raises the 2025 consolidated ratio by a third (2.92 against 2.18).

\subsection{The level: how much inflation would hold debt stable today?}\label{sec:level}

The combined metric is a sensitivity: inflation per point of a permanent rate shock. A policymaker's question is a level: at today's deficits and today's rates, how much inflation would stop debt/GDP from rising? The same two clocks answer it. Hold interest-bearing consolidated debt $b$ and the zero-interest base $z$ (currency, and reserves before 2008) constant relative to GDP for $H$ years. The debt-stabilizing primary balance at horizon $h$ is $(\bar r(h)-g)\,b-g\,z$, where the average rate moves from today's $\bar r_0$ toward the marginal cost along the clock, $\bar r(h)=\bar r_0+(r-\bar r_0)P(h)$, and the growth of zero-interest money, $g\,z$, is seigniorage that finances part of the deficit. The shortfall against the actual primary balance $s(h)$ must be covered by erosion $\Delta\pi\,b\,E(h)$:
\begin{equation}
\Delta\pi^{level}(H)=\frac{\int_0^H\big[(\bar r(h)-g)\,b-g\,z-s(h)\big]\,dh}{b\int_0^H E(h)\,dh}.
\end{equation}
For end-2025, $s(h)$ is CBO's August 2026 baseline primary balance for fiscal years 2026--2035. For 2007 and 2019 we hold the actual primary balance (excluding Fed remittances) constant. The average rate $\bar r_0$ prices each privately held security at its own rate, calibrated to the Treasury's official average rate on marketable debt; TIPS add 2\% expected inflation, reserves earn the interest rate on reserves, and reverse repos the ON RRP rate.

\begin{table}[htbp]
\centering
\caption{Inflation needed to hold consolidated debt/GDP constant}\label{tab:level}
\small
\begin{tabularx}{\textwidth}{>{\raggedright\arraybackslash}p{0.3\textwidth}>{\centering\arraybackslash}X>{\centering\arraybackslash}X>{\centering\arraybackslash}X}
\toprule
 & End-2007 & End-2019 & End-2025 (CBO baseline) \\
\midrule
Interest-bearing consolidated debt, \% of GDP & 26 & 75 & 95 \\
Zero-interest base (currency; reserves before 2008), \% of GDP & 5.8 & 8.4 & 7.9 \\
Average rate on the stock $\bar r_0$, \% & 4.9 & 2.3 & 3.6 \\
Marginal cost r / growth g, \% & 4.6 / 3.1 & 2.1 / 3.9 & 4.2 / 3.8 \\
Primary balance, \% of GDP (10-year mean) & +0.97 & $-$2.38 & $-$1.41 \\
Debt-stabilizing primary balance, \% of GDP (year 0 $\to$ year 10) & +0.29 $\to$ +0.22 & $-$1.52 $\to$ $-$1.64 & $-$0.53 $\to$ +0.03 \\
Shortfall, \% of GDP (10-year mean) & $-$0.74 & 0.76 & 1.32 \\
Erosion per pp of inflation, \$bn a year (10-year mean) & 18 & 58 & 98 \\
Inflation to hold debt/GDP, pp a year, H = 10 (H = 5) & none & 2.8 (2.1) & 4.2 (3.0) \\
Same, discounted at r $-$ g & none & 2.9 & 4.1 \\
\bottomrule
\end{tabularx}
\end{table}


At end-2025, holding debt/GDP constant through inflation alone would take about 4pp of extra inflation a year for a decade (3.0 over five years), against 2.8pp in 2019 and nothing in 2007, when the primary surplus exceeded the debt-stabilizing level. The primary shortfall averages 1.3\% of GDP, about \$400bn a year. Each point of sustained surprise inflation erodes about \$100bn a year on average, more at first and less as the debt reprices, which is why the requirement rises with the horizon. The number does not use $\hat\varphi$, since the actual and projected primary balances already embed the fiscal response. It is an order of magnitude, not a forecast: it holds the primary balance fixed in real terms, and bracket creep and lags in the indexation of spending would reduce it; it holds currency demand and real rates fixed, and a flight from currency or a rise in the inflation risk premium would raise it.
